\subsection{Анализ существующих подходов и средств индексации блокчейн данных}

Существующие подходы к работе с блокчейн данными можно разделить на несколько
групп. Первый подход заключается в прямом обращении прикладной системы к RPC
интерфейсу узла. Его преимуществом является минимальный порог внедрения,
однако при росте объема запросов он приводит к высокой нагрузке на узел,
усложняет реализацию прикладной логики и не обеспечивает готовых индексов для
быстрых выборок по адресам, транзакциям и временным диапазонам.

Второй подход основан на использовании готовых внешних блокчейн-эксплореров и
публичных API. Такое решение снижает затраты на собственную инфраструктуру, но
создает зависимость от стороннего поставщика данных, ограничивает контроль над
полнотой и актуальностью информации, затрудняет реализацию специфических
требований к безопасности и не позволяет формировать собственную модель
хранения, ориентированную на прикладные сценарии.

Третий подход предполагает создание собственного индексатора, который получает
данные от доверенного узла, нормализует их, сохраняет в специализированном
хранилище и предоставляет прикладной \gls{api}. Для задач ВКР именно этот
подход является наиболее обоснованным, поскольку позволяет:
\begin{enumerate}
    \item контролировать полноту и достоверность источника данных;
    \item проектировать модель хранения под конкретные сценарии выборок;
    \item учитывать особенности mempool и reorg;
    \item реализовать управляемый жизненный цикл заданий индексирования;
    \item обеспечить независимость прикладных сервисов от прямого доступа к узлу.
\end{enumerate}

Для сопоставления перечисленных подходов целесообразно учитывать следующие
критерии: контроль над источником данных, полнота исторической информации,
устойчивость к reorg, пригодность для адресной аналитики, управляемость
жизненного цикла обработки и эксплуатационная независимость от внешних
поставщиков. С этих позиций прямой доступ к RPC-интерфейсу целесообразен
только для ограниченного круга технических задач, внешние API удобны при
быстрой интеграции, а собственный индексатор является наиболее подходящим
вариантом для построения устойчивого прикладного контура данных.

Обобщенное сопоставление рассмотренных подходов приведено в
таблице~\ref{tab:indexing-approaches}.

\begin{table}
    \caption{Сравнение подходов к получению и предоставлению блокчейн данных}
    \begin{tabular}{|p{0.26\textwidth}|p{0.2\textwidth}|p{0.2\textwidth}|p{0.2\textwidth}|}\hline
    Критерий & Прямой RPC-доступ & Внешние API & Собственный индексатор \\ \hline
    Контроль над данными & высокий & низкий & высокий \\ \hline
    Полнота исторических представлений & ограниченная & зависит от провайдера & высокая \\ \hline
    Устойчивость к reorg и mempool-состояниям & требует ручной реализации & зависит от провайдера & контролируется в системе \\ \hline
    Пригодность для адресной аналитики & низкая & средняя & высокая \\ \hline
    Эксплуатационная независимость & высокая & низкая & высокая \\ \hline
    \end{tabular}
    \label{tab:indexing-approaches}
\end{table}

В рамках рассматриваемого проекта дополнительно учитываются требования к
архитектуре программной системы. В качестве архитектурного стандарта выбран
модульный монолит со stateless-характером исполнения, что позволяет, с одной
стороны, сохранить целостность прикладной логики в пределах одного приложения,
а с другой стороны,
явно разделить ответственность между модулями индексации, хранения, API,
мониторинга и администрирования. Такой подход соответствует текущему масштабу
разрабатываемой системы и одновременно оставляет возможности для дальнейшего
расширения решения.

Следовательно, для предметной области индексации блокчейн данных оптимальным
является подход, при котором система строится как специализированный программный
слой между узлом Bitcoin и внешними потребителями данных.
